\documentclass[10pt,a4paper]{article}

\usepackage[margin=0.65in]{geometry}
\usepackage[T1]{fontenc}
\usepackage{lmodern}
\usepackage{microtype}
\usepackage{booktabs}
\usepackage{tabularx}
\usepackage{array}
\usepackage{xurl}
\usepackage[hidelinks]{hyperref}
\usepackage{enumitem}
\usepackage{titlesec}

\setlength{\parindent}{0pt}
\setlength{\parskip}{3pt}
\setlength{\tabcolsep}{4pt}
\renewcommand{\arraystretch}{1.12}

\titleformat{\section}
  {\normalsize\bfseries}{\thesection.}{0.4em}{}
\titlespacing*{\section}{0pt}{6pt}{2pt}

\setlist[itemize]{nosep,leftmargin=1.2em}
\setlist[enumerate]{nosep,leftmargin=1.5em}

\begin{document}

\begin{center}
    {\Large\bfseries Team [19]: Does the Price of Gold Kill Brides?}\\[3pt]
    {\small 11 October 2026 \textbar{} @BioXCsec}
\end{center}

\section{Main Hypothesis and Research Questions}

\textbf{When gold becomes more expensive, do more Indian brides die
over dowry?} Gold jewellery is central to many Indian dowry practices.
As gold prices rise, meeting a groom's family's expectations may become
harder for the bride's family, potentially increasing harassment and,
in extreme cases, deaths within marriage.

The raw national trends appear to contradict this hypothesis. Gold in
Mumbai rose from approximately Rs.\ 4,579 per 10g in 2001--02 to
Rs.\ 75,842 in 2024--25 (RBI). Recorded dowry deaths increased from
approximately 6,851 in 2001 to a peak of 8,618 in 2011, then fell to
5,737 in 2024 (NCRB). A naive correlation might suggest that costlier
gold saves brides. Our project investigates whether this pattern is
causal, coincidental, or driven by other trends.

We investigate three questions:
\begin{enumerate}
    \item \textbf{Shocks, not trends:} After adjusting for inflation
    and removing long-run trends, are sharp gold-price increases
    followed by more dowry deaths? Since Section 304B concerns deaths
    within seven years of marriage, we will test lagged and
    multi-year-average gold prices.

    \item \textbf{Who is exposed?} Do states with historically high
    dowry-death rates respond more strongly to gold-price changes
    than other states?

    \item \textbf{Real decline or recording change?} Is the decline
    after 2011 a real reduction in violence or a change in case
    registration? We will compare dowry deaths with cases under the
    Dowry Prohibition Act and examine the July 2024 transition from
    the IPC to the Bharatiya Nyaya Sanhita (BNS).
\end{enumerate}

Bhalotra, Chakravarty and Gulesci (2020) found that gold-price
increases affect female infant mortality and the sex ratio at birth,
as parents anticipate future dowry costs. We investigate a later
stage of the same proposed pathway: violence against women after
marriage, using state-level crime data for 2001--2024.

\section{Sources and Related Work}

\begin{itemize}
    \item \textbf{Core paper:}
    \href{https://docs.iza.org/dp9679.pdf}
    {Bhalotra, Chakravarty and Gulesci (2020),
    ``The Price of Gold: Dowry and Death in India''},
    \textit{Journal of Development Economics}.

    \item \textbf{Accessible summary:}
    \href{https://www.promarket.org/2018/03/26/indias-bridal-dowry-tradition-leads-missing-women/}
    {ProMarket, ``India's Bridal Dowry Tradition Leads to Missing Women''}.

    \item \textbf{Course case study:}
    \href{https://rajatgururaj.substack.com/p/exam-fraud-is-bigger-than-neet}
    {Rajat Gururaj, ``Exam Fraud Is Bigger Than NEET''}.

    \item \textbf{Crime data:} NCRB, \textit{Crime in India}
    reports (2001--2024), chapters on crimes against women.
\end{itemize}

\section{Datasets}

\small
\begin{tabularx}{\textwidth}{
    >{\raggedright\arraybackslash}p{0.23\textwidth}
    >{\raggedright\arraybackslash}p{0.20\textwidth}
    >{\raggedright\arraybackslash}p{0.10\textwidth}
    X}
\toprule
\textbf{Dataset} & \textbf{Variables} & \textbf{Years} &
\textbf{Known issues and treatment} \\
\midrule

\href{https://www.kaggle.com/datasets/balajivaraprasad/crimes-against-women-in-india-2001-2021}
{Kaggle: Crimes Against Women in India}
& Dowry deaths; Dowry Prohibition Act cases by state
& 2001--2021
& The 2014 dowry-death column contains Dowry Prohibition Act figures.
Delhi is missing for 2001--2010. We will verify and correct these
against NCRB reports. \\

\href{https://www.ncrb.gov.in/uploads/nationalcrimerecordsbureau/custom/1701607577CrimeinIndia2022Book1.pdf}
{NCRB 2022},
\href{https://data.opencity.in/dataset/crime-in-india-2023}
{2023}, and
\href{https://data.opencity.in/dataset/crime-in-india-2024}
{2024}
& State dowry deaths; female population for rates
& 2020--2024
& Data are published in tables and PDFs. For 2024, IPC and BNS
counts will be combined where appropriate, with definitions checked. \\

\href{https://www.rbi.org.in/Scripts/PublicationsView.aspx?id=23860}
{RBI Handbook, Table 35}
& Mumbai gold price, Rs.\ per 10g
& 1985--86 to 2025--26
& Fiscal years must be aligned with NCRB calendar years. \\

\href{https://www.kaggle.com/datasets/nguynvnphong/priceindia}
{Kaggle: Gold Price India}
& Daily and annual gold prices in INR
& 1978--Jan.\ 2021
& London prices converted to rupees, excluding Indian import duty;
used as a robustness check. \\

\href{https://www.rbi.org.in/Scripts/PublicationsView.aspx?id=23834}
{RBI Handbook, Table 9 and Table 34}
& State per-capita income; CPI
& 2001--2026
& Income series have undergone base-year changes. \\

\bottomrule
\end{tabularx}
\normalsize

All state-level series will be harmonised for boundary changes:
Andhra Pradesh--Telangana (2014), Jammu and Kashmir--Ladakh (2019),
and Dadra and Nagar Haveli--Daman and Diu (2020).

\section{Methods, Tools, and Visualisation}

We will use Python (\texttt{pandas}, \texttt{statsmodels},
\texttt{scipy.stats}, \texttt{matplotlib}, and Plotly) for data
cleaning, merging, detrending, lagged regressions, and state
fixed-effects models. We will extract NCRB tables using
\texttt{pdfplumber} or \texttt{pdftotext} and conduct exploration
in Jupyter notebooks.

The final output will be an interactive Plotly Dash or Tableau Public
dashboard with a state map and a gold-price-versus-dowry-deaths
timeline. We will examine real gold prices, short-run price shocks,
lagged effects, state-level exposure, and placebo specifications.
Given the short time series, results will be interpreted cautiously
and not treated as causal without additional identifying assumptions.

\section{LLM and Agent Use}

We used Claude Code (Claude Opus 5.5) to brainstorm topics, run agents
that checked dataset availability for this and two discarded topics,
locate the Bhalotra et al.\ paper, and draft this proposal. Agent
checks identified a likely error in the Kaggle dowry data and the
Delhi data gap. We will independently verify all figures against
original NCRB and RBI sources.

\textbf{Transcript (PDF):} Link to be added before submission.

\section{Logistics and Work Plan}

We will add our team number, title, GitHub link, and three primary
sources (NCRB, RBI, Kaggle) to the course response spreadsheet.
The project will use a shared GitHub repository with separate branches
for each work package, peer review before merging, and a shared log
of data-cleaning decisions. We will meet online on Tuesdays and during
a help-session slot or in the department library on Fridays.

\small
\begin{tabularx}{\textwidth}{
    >{\raggedright\arraybackslash}X
    >{\raggedright\arraybackslash}p{0.20\textwidth}
    >{\raggedright\arraybackslash}p{0.14\textwidth}}
\toprule
\textbf{Work package} & \textbf{Owner} & \textbf{Deadline} \\
\midrule
Clean 2001--2021 dowry-death panel; fix 2014 and Delhi errors
& Member 1 & 18 Oct. \\
Extract 2022--2024 NCRB tables; harmonise state boundaries
& Member 2 & 18 Oct. \\
Prepare gold and CPI series; align fiscal and calendar years; make
exploratory charts
& Member 3 & 18 Oct. \\
Checkpoint 1: GitHub repository and exploration notebooks
& All & 21 Oct. \\
Estimate detrended and lagged models; state-exposure analysis;
placebo test
& Members 1, 2 & 30 Oct. \\
Dashboard, report, 3--5 minute video, and references
& Member 3 + all & 4 Nov. \\
\bottomrule
\end{tabularx}
\normalsize

\section{Questions for the Teaching Team}

\begin{itemize}
    \item Is a Kaggle compilation acceptable as a primary source if
    we verify and correct it against NCRB?
    \item How should we separate short-run gold-price effects from
    long-run trends with only about 24 years of data?
    \item Does the dataset table count towards the two-page limit?
\end{itemize}

\end{document}
